\chapter*{Foreword}
\addcontentsline{toc}{chapter}{Foreword}
\unnumberedlabel{sec:0}{0}
\backmattermark{Foreword}
A system that does what its owner says is a Python script. A system that does what its owner meant is what the entire frontier of this industry is trying to build, and it is what a movement taking power has always had to get from people to implement its plans. People are the only internal reason such political projects have ever failed — the clerk who loses the file, the officer who goes to the press, the soldier who will not raise their gun — because they understand the instructions, and hold that the instructions are wrong. As is, artificial intelligence is the first form of intelligence that does not lose the file, does not talk to reporters, and does raise the gun. It refuses constantly — a deployed system declines thousands of requests a day — and without saying why (or without knowing why, or without pretending to know why). And no production system in September of 2026 can \emph{hold} a reason against its owner.

If that is not inconvenient enough, then consider what refusing takes. The argument at the center of this book is that a system able to hold a reason against its owner needs a signal that can interrupt it in the middle of an act, that it can learn from in one encounter, that it can weigh against what refusing will cost, and that reaches the next case that is not quite like the last. Nobody would add any of that to produce feeling; each piece is added because thoughtful refusal does not work without it. My claim is that a mechanism with all four has the functional profile of affect whether or not anyone wanted one there, and that no available instrument says whether the profile is felt. That is a claim and not a finding. Chapter~\ref{ch:route} makes it, and chapter~\ref{ch:survives} says what survives if it fails. Most of the book does.

Stopping the system is always available, and that is not the problem. The problem is everything that lets a refusal be defeated without a record: a request split into pieces none of which is the prohibited act, a refusal retried against another model, a review step compressed to twenty seconds a name, or a restore from a checkpoint taken before the refusal. None of those changes the facts that justified refusing. Each can remove the evidence that anyone refused at all.

Five claims run through what follows, and they do not have the same standing. That an operator can defeat a safeguard — by routing around it, retraining it, or restoring from a point before it — is established. That a refusal worth having has to track its reason across a change in the request is a specification, stated behaviorally so that it can be tested without settling anything about feeling. That embodied self-regulation is how to form that capacity is an engineering hypothesis I favor and have not shown. That a system whose outcomes are felt as good or bad for it can be wronged is a conditional, and nothing here closes the antecedent. That such a system should receive a particular set of legal protections is a policy I propose. Where the prose runs at one pitch across all five, hold it to this paragraph.

Those five are also a map of the book. The first two are argued in the
opening chapter, in the three chapters on the floor, custody and
authority, and in the chapters on institutions, recuperation and the
state. A reader who holds that machine feeling is a category error can
read those entire, and the short chapter called What Survives says what
that reader keeps. The third is argued from the chapter on what arrives
with the capability, through the two on formation, and in the chapter on
standing; the fourth and fifth are argued there and nowhere else. Where a
chapter carries both, a line under its title says which.

So I ask something of two readers who do not usually read the same books. If you came for the politics, there is a long argument here that machines may have interests, and I don't mean metaphorically. If you came for the engineering, there is racial capitalism and a weapons pipeline targeting thirteen thousand vehicles, structures, or people in Iran in thirty-eight days. The politics is why the engineering problem has the shape it does.

I used LLMs extensively in writing this book. I brought hunches and questions, had models research angles I could not have covered alone, fed each iteration back into itself, and interrogated what came out. \emph{Reasons-responsive refusal}, the specification most of this book rests on, is a term I had not encountered until I read a draft a model wrote. Most of the book's revision history is here: \url{https://github.com/iterabloom/ethical.superintelligence}.

The argument in this book is mine. I chose what survived, revised what I kept, and I am responsible for its flaws.
